\begin{theorem}[Conductor nonvanishing for the standard family]
\hypertarget{tag-03IX}{}\label{wbII:thm:candidate-conductor}
\StatementRecord{wbII:thm:candidate-conductor}{proved}{03IX}{\SeriesRef{wbII:eq:W:5.3}, \SeriesRef{wbII:prop:conductor}, \SeriesRef{wbII:thm:cartier-duality}, \SeriesRef{wbII:eq:W:5.20}, \SeriesRef{wbII:lem:curve-flatness}, \SeriesRef{wbII:prop:analytic-base-change}; \cite[Theorem~9.2.3, 2010 text, p.~410]{CoxLittleSchenck2011}; \cite[\S12, Theorem~1, p.~19]{Serre1956GAGA}}

Assume:
\begin{enumerate}[label=(H\arabic*)]
\item $X$ is a connected compact smooth complex threefold and
$f:X\to B$ is a nonconstant holomorphic map to a smooth compact curve;
\item $W=f^{-1}(p_0)$ is the full scheme-theoretic fibre, is reduced, and
has the normal-crossings local models \eqref{wbII:eq:W:5.3};
\item the normalization $\eta:\widetilde W\to W$ is the degree-six del
Pezzo surface $\mathrm{dP}_6$;
\item $H^2(X;\mathbb Z)=0$.
\end{enumerate}
Then, for every $L\in\operatorname{Pic}(X)$,
\[
H^2(W,(T_X\otimes L)|_W)\ne0
\quad\text{and}\quad
R^2f_*(T_X\otimes L)\ne0.
\tag{W.45}\hypertarget{tag-03HI}{}\label{wbII:eq:W:5.18}
\]
\end{theorem}



\begin{proof}

\begin{enumerate}[label=\arabic*.,leftmargin=*]
\item Fix $L$ and put
\[
\mathcal F_L=T_X\otimes L,\qquad
A_L=(L^\vee\otimes\omega_X)|_W.
\]
Hypothesis (H4) gives $c_1(L^\vee\otimes\omega_X)=0$.  Pullback therefore
gives $c_1(\eta^*A_L)=0$. The surface $\mathrm{dP}_6$ is the blow-up of
$\mathbb P^2$ at three noncollinear points. A projective change of coordinates
moves these points to the three coordinate points, so this is the smooth
complete toric surface obtained by subdividing the three two-dimensional
cones of the fan of $\mathbb P^2$. Its fan has support $\mathbb R^2$,
and the zero Cartier divisor is nef. Demazure vanishing therefore gives
$H^1(\mathrm{dP}_6,\mathcal O)=0$ in algebraic cohomology
\cite[Theorem~9.2.3, 2010 text, p.~410]{CoxLittleSchenck2011}.
Since this surface is projective and $\mathcal O$ is coherent,
Serre's comparison theorem identifies that group with analytic cohomology
\cite[\S12, Theorem~1, p.~19]{Serre1956GAGA}.
In the exponential sequence, the kernel
of
\[
c_1:\operatorname{Pic}(\mathrm{dP}_6)\longrightarrow
H^2(\mathrm{dP}_6;\mathbb Z)
\]
is the image of $H^1(\mathrm{dP}_6,\mathcal O)$ and is therefore zero.
Consequently
\[
\eta^*A_L\simeq\mathcal O_{\widetilde W}.
\tag{W.46}\hypertarget{tag-03HJ}{}\label{wbII:eq:W:5.19}
\]
\Cref{wbII:prop:conductor} gives
\[
H^0(W,\Omega^1_X|_W\otimes A_L)\ne0.
\tag{W.47}\hypertarget{tag-03HK}{}\label{wbII:eq:W:5.20}
\]

\item Because $W=f^*(p_0)$, one has
$\mathcal O_X(W)=f^*\mathcal O_B(p_0)$, whose restriction to the fibre is
trivial.  Adjunction therefore gives $\omega_W\simeq\omega_X|_W$.
Apply \cref{wbII:thm:cartier-duality} to $E=\mathcal F_L$:
\[
\begin{split}
H^2(W,\mathcal F_L|_W)^\vee
&\simeq H^0(W,(\mathcal F_L|_W)^\vee\otimes\omega_W)\\
&=H^0(W,\Omega^1_X|_W\otimes A_L).
\end{split}
\]
By \eqref{wbII:eq:W:5.20}, the dual space is nonzero, giving the first nonvanishing in \eqref{wbII:eq:W:5.18}.
The same duality theorem gives
$H^3(W,\mathcal F_L|_W)=0$.

\item Since $X$ is compact, $f$ is proper.  By
\cref{wbII:lem:curve-flatness}, $f$ is flat and the vector bundle
$\mathcal F_L$ is $f$-flat.  Thus every hypothesis of
\cref{wbII:prop:analytic-base-change} is satisfied at $p_0$, and
\[
(R^2f_*\mathcal F_L)_{p_0}\otimes\mathbb C(p_0)
\simeq H^2(W,\mathcal F_L|_W)\ne0.
\]
The stalk, and hence the coherent direct-image sheaf in
\eqref{wbII:eq:W:5.18}, is nonzero.

\end{enumerate}
\end{proof}